\documentclass[12pt]{article}

\usepackage[margin=1in]{geometry}
\usepackage{setspace}
\usepackage[round,authoryear]{natbib}
\usepackage{hyperref}

\title{Overdrawing on Experience: Panpsychism and the Limits of First-Person Warrant}
\author{}
\date{}

\begin{document}

\maketitle
\doublespacing

\begin{abstract}
Panpsychism is often motivated by an appeal to the epistemic privilege of consciousness. Physical science seems to describe matter in structural, relational, dispositional, or dynamic terms, while consciousness is known directly from the first-person point of view. This paper argues that this appeal supports less than panpsychism requires. First-person experience establishes consciousness as a datum and discloses ordinary conscious life as structured, content-involving, temporal, affective, embodied, attentive, and often self-involving. It does not, by itself, disclose bare subjectivity, micro-experience, proto-experience, cosmic subjectivity, hidden phenomenal fields, or relations of phenomenal inheritance. Those posits may still be defended, but they must be defended as theoretical posits and judged by their explanatory work. They cannot inherit the special warrant of ordinary first-person consciousness merely by being described as phenomenal or experiential. The conclusion is not that panpsychism is false. It is that one prominent first-person route to panpsychism overdraws on the very datum that gives the view much of its appeal.
\end{abstract}

\section{Introduction: Panpsychism's Appeal and Overreach}

Panpsychism derives much of its contemporary appeal from an epistemic asymmetry. Consciousness, unlike many entities posited by physical science, is known from the first-person point of view. I do not infer that I am in pain in the way I infer the presence of an electron, a gene, or a gravitational wave. I undergo the pain. I live the experience from within. By contrast, physical science often appears to describe matter in terms of structure, relation, disposition, and causal role. It tells us what matter does, but not, at least on one influential line of thought, what matter is intrinsically. If consciousness is the one reality known from the inside, then perhaps consciousness, or something consciousness-like, belongs in our account of fundamental reality.

This thought gives panpsychism its distinctive force. It promises to avoid the apparent miracle of consciousness emerging from an entirely non-conscious base by locating mentality, proto-mentality, or phenomenal character in the nature of matter itself. In its more sophisticated forms, the view is not the crude claim that electrons have tiny human-like minds. Strawson's realistic monism, for example, argues that a serious physicalism must take experience as part of the physical world and cannot treat experiential reality as an inexplicable late arrival \citep{strawson2006}. Goff's Russellian monist defense places panpsychism in a broader argument about the categorical nature of matter and the limits of structural physics \citep{goff2017cfr}.

The appeal to first-person experience also imposes a constraint. If first-person experience is the source of panpsychism's special warrant, then the view may claim that warrant only for what first-person experience actually discloses. What first-person experience discloses is not micro-experience, proto-experience, cosmic subjectivity, hidden phenomenal fields, or relations of phenomenal inheritance between universal and local minds. It discloses ordinary, structured, content-involving consciousness: sensation, affect, thought, attention, temporal flow, embodiment, memory, and self-relation.

This paper argues that panpsychism often overdraws on the evidential authority of first-person experience. The argument is limited. I do not claim that panpsychism is impossible. I do not claim that hidden entities are illegitimate. I do not claim that physics must be able to explain consciousness in purely structural or functional terms. The claim is narrower: first-person experience gives panpsychism a datum, not a metaphysics. It establishes that consciousness exists and discloses certain features of ordinary conscious life. It does not establish that consciousness is fundamental, ubiquitous, microphysical, cosmic, or hidden in the categorical basis of physical structure.

The rest of the paper develops this point in stages. Section 2 fixes the target and brackets arguments that are not being challenged here. Section 3 states the introspective warrant constraint. Section 4 argues that first-person experience discloses structured consciousness rather than bare subjectivity. Section 5 applies the constraint to micro-experience, proto-experience, and cosmic subjectivity. Section 6 explains why Russellian structure/dynamics arguments do not by themselves close the gap. Section 7 distinguishes austere micropsychist views from more elaborate cosmopsychist machinery. Section 8 turns to explanatory burden, using dreamless sleep and temporal grain as conceptual illustrations. Section 9 concludes that panpsychism's first-person motivation is weaker than often advertised.

\section{Scope and Bracketing}

The argument should not be read as an attack on every possible route to panpsychism. Panpsychists may appeal to the alleged impossibility of radical emergence, to continuity in nature, to theoretical simplicity, to dissatisfaction with standard physicalism, or to the explanatory gap. They may also defend a Russellian monist thesis according to which physics describes only structure and dynamics, leaving the intrinsic nature of matter unspecified. Those arguments raise large issues that cannot be settled by the present paper.

My target is narrower. I am concerned with one prominent motivation for panpsychism: the claim that consciousness has a special epistemic status because it is known directly from within. This motivation matters because it gives panpsychism a distinctive rhetorical and philosophical force. The panpsychist does not merely add a speculative property to the world. The panpsychist appears to begin from the one thing we know most intimately, namely experience itself.

That route is legitimate as far as it goes. But it does not go as far as panpsychism often needs it to go. First-person experience may establish the reality of consciousness. It may also disclose features of conscious life that any adequate theory should respect. It does not automatically support claims about the distribution, fundamentality, or cosmic structure of consciousness.

This scope restriction is important because a panpsychist could otherwise object that the paper confuses two questions. One question is whether consciousness exists and is known in a distinctive first-person way. A second question is whether consciousness, or something related to it, belongs to the fundamental nature of matter. The first question can be answered from the first-person point of view. The second cannot. It requires a metaphysical argument.

Nor do I ignore the Russellian structure/dynamics argument entirely. I bracket the full case for Russellian monism, but I address one step within it. Even if physics gives only structure and dynamics, it does not follow that the categorical nature behind that structure is phenomenal or proto-phenomenal. The structure/dynamics argument may create a vacancy. It does not show that consciousness fills it.

\section{The Introspective Warrant Constraint}

Panpsychism often draws force from a claim about the special epistemic status of consciousness. Consciousness is not known in the way electrons, genes, or black holes are known. I do not posit my current pain as a theoretical entity. I undergo it. Conscious experience is given from the first-person point of view.

This appeal is legitimate, but it must be disciplined. The fact that consciousness is known directly does not mean that every entity described as phenomenal, proto-phenomenal, or experiential inherits the warrant of first-person experience. First-person warrant is not a transferable credit. It attaches to what is actually given in experience.

Call this the introspective warrant constraint:

\begin{quote}
A theory that motivates itself by appeal to first-person experience may claim first-person warrant only for what first-person experience actually discloses.
\end{quote}

The constraint is modest. It does not say that only directly experienced entities exist. It does not prohibit theoretical posits. Physics, biology, and psychology all legitimately posit entities and structures not given in ordinary experience. But those posits are warranted by explanatory, predictive, or unificatory work. They do not inherit the authority of first-person acquaintance. If panpsychism posits micro-experience, proto-experience, cosmic subjectivity, hidden phenomenal fields, or relations of phenomenal inheritance, those posits may be defended as theoretical entities. They may not be treated as if they were directly warranted by the first-person datum of consciousness.

The point is clearest if we distinguish three steps. First, there is the datum: consciousness exists. There is something it is like to undergo pain, see red, hear a sound, remember a face, or feel time passing. Second, there is the characterization of the datum: first-person experience presents consciousness as structured, determinate, temporal, affective, embodied, attentive, and often self-involving. Third, there is metaphysical inflation: the claim that consciousness, proto-consciousness, or subjectivity belongs to the fundamental nature of matter, exists at the microphysical level, or characterizes the universe as a whole.

The introspective warrant constraint allows the first step. It also allows the second. But it blocks any automatic passage to the third. First-person experience warrants the existence of ordinary consciousness and discloses something about its character. It does not, by itself, disclose consciousness as fundamental, ubiquitous, microphysical, cosmic, or hidden in the categorical basis of physical reality.

This constraint matters because panpsychism often presents itself as unusually faithful to consciousness. Its promise is to take experience seriously rather than explaining it away. But taking experience seriously also means taking seriously the limits of what experience gives. First-person experience gives us conscious life as lived: sensations, affects, thoughts, moods, temporal flow, attention, memory, embodiment, and self-relation. It does not give us bare subjectivity detached from all phenomenal character, nor micro-experiences hidden in particles, nor proto-experiential categorical properties, nor a unified cosmic experiencer.

\section{What First-Person Experience Discloses}

The first-person datum is sometimes described as bare subjectivity, pure awareness, or mere what-it-is-likeness. That description is tempting because it seems to isolate the minimal feature common to all conscious states. Yet it risks abstracting away precisely the features that make consciousness intelligible as a datum.

When we attend to consciousness, we do not find subjectivity as an empty item standing apart from all character. We find pain, color, warmth, pressure, thought, anticipation, anxiety, boredom, temporal flow, bodily orientation, and degrees of attention. Even when experience seems objectless, it has some character. A state of blankness is not the same as a state of calm. Silence is not the same as visual darkness. Bodily absence is not the same as bodily comfort. The relevant point is not that every conscious state represents an external object. It is that every conscious state we can identify has some phenomenal structure.

This creates a dilemma for appeals to fundamental subjectivity. If fundamental consciousness has phenomenal character, then it already has structure, and that structure calls for explanation. It is not a simple primitive known by introspection alone. If fundamental consciousness has no phenomenal character, then it is not something first-person experience discloses. It is an abstraction from experience, not an item found in experience.

The dilemma is especially important for views that retreat from rich consciousness to minimal subjectivity. The move is understandable. Few panpsychists want to say that electrons have visual fields, memories, anxieties, or personal identities. So the relevant posit is often thinned out. The fundamental level is said to contain minimal phenomenal character, bare subjectivity, or proto-experiential nature. But thinning the posit changes its epistemic status. The thinner it becomes, the less it resembles the consciousness that first-person experience actually gives.

This does not show that minimal subjectivity is incoherent. It shows that it cannot simply borrow the warrant of ordinary experience. The familiar datum is structured consciousness. A stripped-down subjectivity with no content, valence, temporality, mineness, intensity, or structure is not given in the same way. If it is introduced, it is introduced by theory.

\section{The Failure of Warrant Transfer}

The introspective warrant constraint becomes most important when panpsychism moves from ordinary consciousness to hidden phenomenal ontology. Contemporary panpsychist and panprotopsychist views may posit micro-experience, proto-experience, proto-phenomenal properties, cosmic subjectivity, hidden phenomenal fields, phenomenal bonding relations, or forms of phenomenal inheritance. These posits are not directly disclosed in ordinary first-person experience.

The problem is not that such posits are impossible. The problem is that they cannot automatically inherit the epistemic privilege of first-person consciousness. The immediacy of pain warrants the pain. It does not warrant a cosmic subject. The directness of visual experience warrants that there is visual experience. It does not warrant the claim that microphysical entities instantiate hidden phenomenal natures.

Panprotopsychism makes the warrant-transfer problem especially clear. The panprotopsychist need not say that fundamental entities are conscious. The claim is that they instantiate proto-phenomenal or proto-experiential properties that can help explain consciousness. But proto-experience is not something we encounter from within. It is introduced as a theoretical posit. Its relation to ordinary consciousness must therefore be argued for, not assumed.

Here the panprotopsychist faces a dilemma. If proto-experience is sufficiently like experience, then it inherits familiar problems about structure, unity, subjecthood, and combination. If it is sufficiently unlike experience, then it loses its claim to first-person warrant. Either way, the appeal to ordinary consciousness does less work than advertised.

The same point applies to phenomenal bonding. Goff's discussion of phenomenal bonding is explicitly designed to address the combination problem, the problem of how many micro-level subjects or phenomenal properties could yield one macro-level subject \citep{goff2017bonding}. That machinery may be part of a serious metaphysical theory. But it is not given in ordinary first-person experience. I am acquainted with pain, color, thought, and temporal flow. I am not acquainted, simply by being conscious, with relations by which micro-subjects combine into macro-subjects.

Panpsychism overdraws on first-person evidence when it treats the epistemic privilege of ordinary consciousness as support for a hidden phenomenal ontology. The hidden ontology may still be defended. But it must earn its warrant as theory.

\section{Structure and Dynamics Do Not Rescue the Inference}

The Russellian monist can press a further argument. Physics, it is said, gives us only structure and dynamics. It describes relations, dispositions, equations, and causal roles. It leaves open the intrinsic or categorical nature of the entities that bear those roles. Consciousness, by contrast, is known from within. Perhaps consciousness reveals what matter is intrinsically, or at least gives us a model for what intrinsic nature could be \citep{goff2017cfr}.

This argument deserves careful treatment. It is not a crude appeal to mystery. It begins from a real limitation in structural description. If a theory tells us only how entities are related or how they behave, one may ask what, if anything, has the structure or grounds the dispositions. The question is intelligible.

But the step from vacancy to phenomenal filling is not licensed by first-person experience alone. From the claim that matter has an unknown intrinsic nature, it does not follow that this intrinsic nature is experiential or proto-experiential. First-person experience discloses ordinary consciousness. It does not disclose the intrinsic nature of electrons, fields, spacetime, or the cosmos. To identify categorical physical nature with phenomenal or proto-phenomenal nature is therefore a metaphysical inference.

This point does not refute Russellian monism. A Russellian panpsychist may have further arguments. The claim may be defended by appeal to simplicity, explanatory gap considerations, anti-emergentist premises, or a broader theory of categorical properties. But once the argument takes that form, it is no longer simply an appeal to first-person immediacy. It is a theoretical proposal about what best fills an explanatory or metaphysical gap.

The distinction matters because first-person evidence has unusual authority only within its proper domain. I can be directly aware of my pain. I cannot be directly aware, by that same act, that the intrinsic nature of all matter is pain-like, experience-like, or proto-phenomenal. The structure/dynamics argument may explain why a panpsychist looks for something beyond structural physics. It does not show that the thing found must be consciousness.

\section{Goff, Strawson, and Degrees of Violation}

The introspective warrant problem applies with different force to different panpsychist views. This variation matters. A fair criticism should not treat every version of panpsychism as if it made the same claim.

Strawson's realistic monism is the austere baseline. Strawson argues that experience is real, that experience is physical if physicalism is true, and that the emergence of experience from the wholly non-experiential would be unintelligible \citep{strawson2006}. This argument does not depend on cartoon particle minds. It is a serious attempt to force physicalism to include experience within the physical.

The present objection still applies, but it applies in a specific way. Strawson may be right that experience is part of reality and cannot be dismissed. The question is whether first-person acquaintance with ordinary experience warrants the stronger claim that experientiality exists at the fundamental level. The introspective warrant constraint says no. It permits the datum of experience. It does not permit an automatic inference from that datum to microphysical experientiality.

Goff's more elaborate Russellian and cosmopsychist machinery makes the problem sharper. His work develops panpsychism within a structure/dynamics framework and addresses the combination problem through proposals such as phenomenal bonding and, in later discussion, cosmopsychism \citep{goff2017cfr,goff2017bonding}. Cosmopsychism is especially important because it reverses the micropsychist direction. Instead of beginning with many micro-subjects and trying to combine them into human subjects, it begins with a cosmic subject and attempts to derive local subjects from it.

This may avoid some versions of the combination problem. But it worsens the introspective warrant problem. Cosmic subjectivity is not merely hidden from third-person science. It is unlike anything first-person experience discloses. Ordinary first-person experience reveals local, finite, perspectival consciousness with limited contents, attention, memory, and temporal organization. It does not reveal a unified cosmic experiencer from which local minds derive.

The more panpsychism elaborates the machinery needed to solve its internal problems, the farther it moves from the first-person datum that gave it plausibility. That does not show that the machinery is false. It shows that it cannot retain the special warrant of direct first-person evidence. It becomes speculative metaphysics.

\section{Explanatory Burden: Sleep and Temporal Grain as Illustrations}

Once panpsychist posits are treated as theoretical rather than first-personal, the question changes. The issue is no longer whether micro-experience, proto-experience, or cosmic subjectivity is possible. The issue is what these posits explain.

An adequate theory of consciousness should help explain why consciousness is unified, why it has contents, why it has temporal grain, why it is affected by sleep, anesthesia, attention, memory, trauma, and brain damage, why some processing is conscious while other processing is not, why there is one subject rather than many, and why experience has valence. Panpsychism does not answer these questions merely by locating experience, proto-experience, or phenomenal character at the fundamental level.

In many cases, the explanatory work still appears to be done by organization and processing. Unity is explained, if at all, by integration. Temporal grain is explained, if at all, by the temporal organization of processing. Memory, attention, reportability, and access depend on system-level architecture. Valence plausibly depends on organism-level regulation, evaluation, and need. The hidden phenomenal layer is not competing with these explanations. It is standing beside them while they do the work.

This complaint belongs near Papineau's broader treatment of the explanatory gap and phenomenal concepts. Papineau argues that the felt gap between physical and phenomenal description may arise from how we think about conscious states, not from the need to posit extra non-physical ontology \citep{papineau2002}. The present argument is not identical to Papineau's physicalism. It does, however, share a pressure point: before adding hidden phenomenal properties to reality, we should ask what explanatory work they perform.

Dreamless sleep is a useful conceptual illustration, provided it is handled carefully. Dreamless sleep does not prove that consciousness is absent. A panpsychist, a Buddhist philosopher, or a theorist of minimal consciousness may argue that some form of experience occurs without later memory. The modest point is that such a claim requires independent warrant. If there is no memory, no report, no phenomenological trace, and no distinctive explanatory payoff, then first-person evidence does not license the posit.

Temporal grain illustrates the same pattern. Human consciousness appears to have thresholds and integration windows. Some stimuli may be processed without entering experience. If the panpsychist responds by assigning hidden experientiality below the threshold of reportable consciousness, that response has not explained conscious temporal structure. It has multiplied hidden phenomenal facts beneath it.

These examples are not empirical proofs against panpsychism. They are illustrations of the warrant problem. Panpsychism repeatedly posits experience or experience-like properties where experience itself gives no clear warrant. That may be permissible as theory. It is not directly licensed by the first-person datum.

The point is strengthened, not weakened, by worries about introspection itself. Schwitzgebel argues that introspection may be unreliable even in ordinary cases where we might expect it to be secure \citep{schwitzgebel2008,schwitzgebel2011}. One need not accept the strongest version of that claim to see its relevance. If introspection is fallible even within its proper domain, then it should be extended beyond that domain only with caution.

\section{Conclusion: The Forfeiture of Special Warrant}

Panpsychism is not thereby refuted. The argument of this paper is narrower. It concerns the warrant supplied by first-person experience. First-person experience establishes that consciousness exists and discloses certain features of ordinary conscious life. It does not establish that consciousness is fundamental, ubiquitous, microphysical, cosmic, proto-phenomenal, or hidden in the categorical basis of physical structure.

This distinction changes the dialectic. The issue is not whether consciousness is real. It is. The issue is whether panpsychism's hidden phenomenal ontology is supported by the same first-person evidence that supports the reality of consciousness. It is not. That ontology may still be defended, but only as theory. It must earn its keep.

The result is a constraint rather than a disproof. Panpsychism may appeal to independent metaphysical arguments, anti-emergentist principles, simplicity considerations, or Russellian claims about intrinsic nature. What it may not do is borrow the authority of ordinary first-person consciousness and spend it on entities or structures that ordinary first-person consciousness never reveals.

Panpsychism overdraws on experience when it treats the epistemic privilege of ordinary consciousness as support for a hidden phenomenal ontology. Those further posits may be possible, but they do not have the warrant panpsychism often implies they have. First-person experience gives panpsychism a datum, not a metaphysics.

\bibliographystyle{plainnat}
\bibliography{overdrawing-on-experience}

\end{document}
